
Length-controlled evaluation metrics have been proposed to debias LLM preference assessments, but whether they preserve the capability ordering that human annotators reveal has not been tested with confound-controlled methods at population scale. This study asks: after explicitly removing the effect of response verbosity, does human preference win rate (win\_rate) still predict length-debiased win rate (LC\_winrate) --- and how strongly? Analyzing the AlpacaEval 2.0 leaderboard across 223 diverse LLMs, a Spearman partial correlation of 0.985 (p = 1.69e-170, one-tailed, bootstrap 95\% CI [0.976, 0.988]) is observed after controlling for avg\_length. This is higher than the 0.94 bivariate correlation previously reported by Dubois et al. [2024] and higher than the bivariate correlation of 0.966 observed in the present sample. In standardized regression, capability explains approximately 4.9 times more LC variance than verbosity (|$\beta _win|$ = 21.34 vs |$\beta _len|$ = 4.37), with a negative verbosity coefficient ($\beta _len = -4.37)$ confirming the directional effect of the length-control correction. Monotonic LC ordering across capability quartiles carries a large effect size ($\epsilon ^2$ = 0.883, KW H = 196.32, p = 2.63e-42). Two independent statistical estimators converge within 1.1 percentage points (Frisch-Waugh-Lovell-inspired verification; $\rho_{resid}$ = 0.9739, delta = 0.011). The alignment gap $\Delta$ = LC\_winrate - win\_rate varies significantly across quartiles (KW p = 5.97e-05) but is non-monotonic and carries only a medium effect size ($\epsilon ^2$ = 0.088), providing a concrete illustration of how dependent variable choice attenuates effect estimates by approximately $10\times$ in capability-alignment studies.

